\documentclass{ximera}

\input{../../../preamble.tex}


\definecolor{fvdnavy}{HTML}{071D43}
\definecolor{fvdcyan}{HTML}{20BFC6}
\definecolor{fvdcyanlight}{HTML}{EFFCFB}
\definecolor{fvdmagenta}{HTML}{F10D69}
\definecolor{fvdmagentalight}{HTML}{FFF1F7}
\definecolor{fvdtext}{HTML}{163044}





%=================================================
% Pedagogische omgevingen
% PDF: tcolorbox
% HTML: learning-box
%=================================================

\newenvironment{voorkennis}
{%
    \ifdefined\HCode
        \HCode{%
<section class="learning-box learning-box--prior">
<div class="learning-box__header">
<span class="learning-box__icon" aria-hidden="true"></span>
<span class="learning-box__title">Voorkennis</span>
</div>
<div class="learning-box__content">
        }%
        \begin{itemize}
    \else
       \begin{tcolorbox}[
    enhanced,
    breakable,
    colback=fvdcyanlight,
    colframe=fvdcyan,
    coltext=fvdtext,
    boxrule=0.5pt,
    leftrule=4pt,
    arc=4mm,
    outer arc=4mm,
    left=7mm,
    right=6mm,
    top=5mm,
    bottom=5mm,
    before skip=6mm,
    after skip=6mm,
    title=Voorkennis,
    coltitle=fvdcyan,
    fonttitle=\bfseries\Large,
    attach title to upper={\par\smallskip},
    boxed title style={
        empty,
        left=0mm,
        right=0mm,
        top=0mm,
        bottom=0mm
    },
    overlay={
        \draw[
            fvdcyan,
            line width=0.5pt
        ]
        ([xshift=42mm,yshift=-13mm]frame.north west)
        --
        ([xshift=-7mm,yshift=-13mm]frame.north east);
    }
]
        \begin{itemize}
    \fi
}
{%
    \end{itemize}
    \ifdefined\HCode
        \HCode{%
</div>
</section>
        }%
    \else
        \end{tcolorbox}
    \fi
}


\newenvironment{leerdoelen}
{%
    \ifdefined\HCode
        \HCode{%
<section class="learning-box learning-box--goals">
<div class="learning-box__header">
<span class="learning-box__icon" aria-hidden="true"></span>
<span class="learning-box__title">Leerdoelen</span>
</div>
<div class="learning-box__content">
        }%
        \begin{itemize}
    \else
\begin{tcolorbox}[
    enhanced,
    breakable,
    colback=fvdmagentalight,
    colframe=fvdmagenta,
    coltext=fvdtext,
    boxrule=0.5pt,
    leftrule=4pt,
    arc=4mm,
    outer arc=4mm,
    left=7mm,
    right=6mm,
    top=5mm,
    bottom=5mm,
    before skip=6mm,
    after skip=6mm,
    title=Leerdoelen,
    coltitle=fvdmagenta,
    fonttitle=\bfseries\Large,
    attach title to upper={\par\smallskip},
    boxed title style={
        empty,
        left=0mm,
        right=0mm,
        top=0mm,
        bottom=0mm
    },
    overlay={
        \draw[
            fvdmagenta,
            line width=0.5pt
        ]
        ([xshift=42mm,yshift=-13mm]frame.north west)
        --
        ([xshift=-7mm,yshift=-13mm]frame.north east);
    }
]
        \begin{itemize}
    \fi
}
{%
    \end{itemize}
    \ifdefined\HCode
        \HCode{%
</div>
</section>
        }%
    \else
        \end{tcolorbox}
    \fi
}



\title{Functieonderzoek}
\author{Ingmar Herreman}

\begin{document}

% FVD_CHAPTER_DATA_START
% Automatisch gegenereerd uit index.tex.
% Niet handmatig aanpassen.
\fvdchapterdata{3}{Afleiden van samengestelde functie}{4}
% FVD_CHAPTER_DATA_END





\begin{abstract}
In de vorige hoofdstukken leerde je afgeleiden berekenen met onder meer
de productregel, quotiëntregel, kettingregel en de regel van l'Hôpital.
In dit hoofdstuk gebruiken we die technieken als gereedschap om het
\textbf{verloop van functies te analyseren}.

De centrale vraag is niet langer alleen ``wat is \(f'(x)\)?'', maar:
\emph{wat vertelt \(f'(x)\) en \(f''(x)\) over de functie en haar grafiek?}

We herhalen kort de stelling van Rolle en de middelwaardestelling van
Lagrange, die je eerder voor veeltermfuncties zag, en passen ze nu ook toe
op andere functietypes. Daarna bouwen we stap voor stap naar volledige
functieonderzoeken en extremumproblemen.
\end{abstract}

\maketitle

\begin{voorkennis}
  \item Je kan het domein, nulpunten en tekenverloop van een functie bepalen.
  \item Je kan limieten berekenen en horizontale, verticale en schuine asymptoten onderzoeken.
  \item Je kan de eerste en tweede afgeleide functie berekenen.
  \item Je kan de productregel, quotiëntregel en kettingregel toepassen.
  \item Je kent het verband tussen het teken van \(f'(x)\) en stijgen of dalen.
  \item Je kent de begrippen lokaal maximum, lokaal minimum en buigpunt.
  \item Je hebt de stelling van Rolle en de middelwaardestelling van Lagrange al gezien voor veeltermfuncties.
\end{voorkennis}

\begin{leerdoelen}
  \item Je kan de voorwaarden van de stelling van Rolle controleren en de stelling toepassen op verschillende functietypes.
  \item Je kan de voorwaarden van de middelwaardestelling van Lagrange controleren en de stelling toepassen en interpreteren.
  \item Je kan met \(f'(x)\) intervallen van stijgen en dalen bepalen.
  \item Je kan stationaire punten onderzoeken en beslissen of er een lokaal extremum optreedt.
  \item Je kan met \(f''(x)\) de kromming van een grafiek onderzoeken en buigpunten bepalen.
  \item Je kan informatie over domein, nulpunten, teken, asymptoten, \(f'\) en \(f''\) samenbrengen in een volledig functieonderzoek.
  \item Je kan vanuit analytische informatie een betrouwbare grafiekschets opbouwen.
  \item Je kan volledige functieonderzoeken uitvoeren voor rationale, irrationale, exponentiële, logaritmische en goniometrische functies.
  \item Je kan extremumproblemen mathematiseren, oplossen en het gevonden extremum in de context interpreteren.
  \item Je kan je keuzes en conclusies wiskundig verantwoorden.
\end{leerdoelen}

% ============================================================
% 1. VAN AFGELEIDE NAAR VERLOOP
% ============================================================

\FVDsection{Van afgeleide naar verloop}

De afgeleide \(f'(x)\) beschrijft de lokale verandering van een functie.
Het teken van \(f'(x)\) vertelt daarom rechtstreeks of de grafiek stijgt
of daalt.

\begin{eigenschap}[Teken van de eerste afgeleide]
Voor een afleidbare functie \(f\) op een interval geldt:

\[
f'(x)>0 \Longrightarrow f \text{ is stijgend},
\qquad
f'(x)<0 \Longrightarrow f \text{ is dalend}.
\]

Een punt waarvoor

\[f'(a)=0\]

heet een \textbf{stationair punt}. Zo'n punt kan een lokaal maximum,
een lokaal minimum of een stationair punt zonder extremum zijn.
\end{eigenschap}

\begin{opmerking}[Horizontale raaklijn is niet hetzelfde als extremum]
Uit \(f'(a)=0\) volgt alleen dat de raaklijn horizontaal is.
Voor een extremum moet het teken van \(f'\) rond \(a\) veranderen.

\[
+\ \to\ - \quad \Longrightarrow \quad \text{lokaal maximum},
\]

\[
-\ \to\ + \quad \Longrightarrow \quad \text{lokaal minimum}.
\]

Als het teken niet verandert, is er geen extremum.
\end{opmerking}

\begin{voorbeeld}[Een rationaal voorbeeld]
Onderzoek het stijgen en dalen van

\[f(x)=\frac{x^3}{2-x^2}.\]

Het domein is

\[\operatorname{dom}f=\mathbb{R}\setminus\{-\sqrt2,\sqrt2\}.\]

De afgeleide is

\[
f'(x)
=
\frac{3x^2(2-x^2)-x^3(-2x)}{(2-x^2)^2}
=
\frac{x^2(6-x^2)}{(2-x^2)^2}.
\]

Omdat de noemer in het domein strikt positief is, wordt het teken van
\(f'\) bepaald door

\[x^2(6-x^2).\]

De kritieke waarden zijn \(x=0\) en \(x=\pm\sqrt6\), terwijl
\(x=\pm\sqrt2\) niet tot het domein behoren.

We krijgen:

\[
\begin{array}{c|ccccccccc}
x & (-\infty,-\sqrt6) & -\sqrt6 &
(-\sqrt6,-\sqrt2) & (-\sqrt2,0) & 0 &
(0,\sqrt2) & (\sqrt2,\sqrt6) & \sqrt6 & (\sqrt6,+\infty)\\
\hline
f'(x) & - & 0 & + & + & 0 & + & + & 0 & -\\
\end{array}
\]

Dus \(f\) is dalend op \((-\infty,-\sqrt6)\), stijgend op de delen van
het domein tussen \(-\sqrt6\) en \(\sqrt6\), en opnieuw dalend op
\((\sqrt6,+\infty)\).

Bij \(x=-\sqrt6\) verandert \(f'\) van \(-\) naar \(+\): daar ligt een
lokaal minimum. Bij \(x=\sqrt6\) verandert \(f'\) van \(+\) naar \(-\):
daar ligt een lokaal maximum. Bij \(x=0\) verandert het teken niet, dus
daar is geen extremum.
\end{voorbeeld}

% ============================================================
% 2. ROLLE
% ============================================================

\FVDsection{Korte herhaling: de stelling van Rolle}

De stelling van Rolle ken je al uit het functieonderzoek van
veeltermfuncties. We herhalen ze kort en leggen de nadruk op de
\textbf{voorwaarden} en op toepassingen buiten de veeltermfuncties.

\begin{eigenschap}[Stelling van Rolle]
Als een functie \(f\)

\begin{enumerate}
  \item continu is op het gesloten interval \([a,b]\),
  \item afleidbaar is op het open interval \((a,b)\),
  \item en \(f(a)=f(b)\),
\end{enumerate}

dan bestaat er minstens één \(c\in(a,b)\) waarvoor

\[\boxed{f'(c)=0.}\]
\end{eigenschap}

Geometrisch betekent dit: als de grafiek op de uiteinden van het interval
dezelfde hoogte heeft, dan bestaat er minstens één punt tussenin waar de
raaklijn horizontaal is.

\begin{voorbeeld}[Rolle voor een goniometrische functie]
Beschouw

\[f(x)=\sin x\]

op \([0,\pi]\).

De functie is continu op \([0,\pi]\), afleidbaar op \((0,\pi)\) en

\[f(0)=0=f(\pi).\]

Rolle garandeert dus minstens één \(c\in(0,\pi)\) waarvoor \(f'(c)=0\).

Omdat

\[f'(x)=\cos x,\]

vinden we

\[\cos c=0 \Longrightarrow c=\frac{\pi}{2}.\]
\end{voorbeeld}

\begin{voorbeeld}[Wanneer Rolle niet toepasbaar is]
Beschouw

\[f(x)=\frac1x\]

op \([-1,1]\).

Hoewel

\[f(-1)=f(1)=1 \quad\text{niet geldt,}\]

is er bovendien nog een fundamenteler probleem: \(f\) is niet continu op
\([-1,1]\), want \(x=0\) ligt niet in het domein.

De stelling van Rolle kan dus niet worden toegepast.
\end{voorbeeld}

\begin{opmerking}[De stelling geeft een bestaan, geen uniek punt]
Rolle zegt dat er \emph{minstens één} geschikte \(c\) bestaat.
Er kunnen er meerdere zijn.

De stelling zegt ook niet dat elk stationair punt een extremum is.
\end{opmerking}

% ============================================================
% 3. LAGRANGE
% ============================================================

\FVDsection{Korte herhaling: de middelwaardestelling van Lagrange}

De middelwaardestelling van Lagrange veralgemeent Rolle.

\begin{eigenschap}[Middelwaardestelling van Lagrange]
Als een functie \(f\)

\begin{enumerate}
  \item continu is op \([a,b]\),
  \item afleidbaar is op \((a,b)\),
\end{enumerate}

dan bestaat er minstens één \(c\in(a,b)\) waarvoor

\[
\boxed{f'(c)=\frac{f(b)-f(a)}{b-a}.}
\]

De rechterkant is de richtingscoëfficiënt van de koorde door
\((a,f(a))\) en \((b,f(b))\).
\end{eigenschap}

Geometrisch zegt de stelling dat er ergens tussen \(a\) en \(b\) een
raaklijn bestaat die evenwijdig is met die koorde.

\begin{voorbeeld}[Lagrange voor een exponentiële functie]
Beschouw

\[f(x)=e^x\]

op \([0,1]\).

De voorwaarden zijn vervuld. Er bestaat dus een \(c\in(0,1)\) waarvoor

\[
f'(c)=\frac{e-1}{1-0}=e-1.
\]

Omdat \(f'(x)=e^x\), volgt

\[e^c=e-1.\]

Dus

\[\boxed{c=\ln(e-1).}\]
\end{voorbeeld}

\begin{voorbeeld}[Interpretatie in fysica]
Een deeltje beweegt volgens een plaatsfunctie \(s(t)\). Tussen
\(t=a\) en \(t=b\) is de gemiddelde snelheid

\[
v_{\mathrm{gem}}
=
\frac{s(b)-s(a)}{b-a}.
\]

Als \(s\) continu is op \([a,b]\) en afleidbaar op \((a,b)\), dan zegt
Lagrange dat er minstens één tijdstip \(c\in(a,b)\) bestaat waarop

\[
v(c)=s'(c)=v_{\mathrm{gem}}.
\]

De stelling verbindt dus een gemiddelde verandering over een interval
met een ogenblikkelijke verandering in minstens één punt.
\end{voorbeeld}

\begin{opmerking}[Rolle als bijzonder geval]
Als \(f(a)=f(b)\), dan is

\[
\frac{f(b)-f(a)}{b-a}=0.
\]

Lagrange geeft dan \(f'(c)=0\). Rolle is dus een bijzonder geval van
de middelwaardestelling van Lagrange.
\end{opmerking}

% ============================================================
% 4. TWEEDE AFGELEIDE
% ============================================================

\FVDsection{De tweede afgeleide en de kromming van de grafiek}

De tweede afgeleide vertelt hoe de eerste afgeleide verandert.

\begin{eigenschap}[Kromming]
Op een interval waar \(f''\) bestaat:

\[
f''(x)>0
\]

betekent dat \(f'\) toeneemt. De hellingen van de raaklijnen worden groter.

\[
f''(x)<0
\]

betekent dat \(f'\) afneemt. De hellingen van de raaklijnen worden kleiner.
\end{eigenschap}

In deze cursus gebruiken we de termen \textbf{hol} en \textbf{bol}
volgens de afgesproken conventie in de cursus. Belangrijker dan de naam
is dat je steeds het teken van \(f''\) koppelt aan de verandering van de
helling.

\begin{definitie}[Buigpunt]
Een punt \(B(a,f(a))\) is een \textbf{buigpunt} wanneer de kromming van
de grafiek verandert bij \(x=a\).

Een noodzakelijke controle is dus dat het teken van \(f''\) aan
weerszijden van \(a\) verandert.
\end{definitie}

\begin{opmerking}[Ook hier is nul worden niet voldoende]
Uit

\[f''(a)=0\]

volgt niet automatisch dat \(B(a,f(a))\) een buigpunt is.

Je moet nagaan of het teken van \(f''\) werkelijk verandert.
\end{opmerking}

\begin{voorbeeld}[Stationair buigpunt]
Neem

\[f(x)=x^3.\]

Dan

\[f'(x)=3x^2,\qquad f''(x)=6x.\]

Bij \(x=0\) geldt zowel \(f'(0)=0\) als \(f''(0)=0\).
De tweede afgeleide verandert van negatief naar positief, dus \((0,0)\)
is een buigpunt. Omdat \(f'(0)=0\), heeft de grafiek daar bovendien een
horizontale raaklijn.

Het punt is dus een \textbf{stationair buigpunt}, geen extremum.
\end{voorbeeld}

% ============================================================
% 5. VOLLEDIG FUNCTIEONDERZOEK
% ============================================================

\FVDsection{Een volledig functieonderzoek}

Een volledig functieonderzoek combineert informatie die je in vorige
hoofdstukken afzonderlijk hebt leren bepalen.

Gebruik onderstaande volgorde als \textbf{controlelijst}, niet als
mechanisch voorschrift. Soms kan een stap worden overgeslagen of is een
andere volgorde efficiënter.

\begin{enumerate}
  \item Bepaal het \textbf{domein}.
  \item Onderzoek \textbf{continuïteit} en eventuele domeinonderbrekingen.
  \item Bepaal \textbf{nulpunten} en snijpunten met de assen.
  \item Onderzoek het \textbf{teken van \(f\)}.
  \item Onderzoek eventuele \textbf{symmetrie} of periodiciteit.
  \item Bepaal relevante \textbf{limieten} en \textbf{asymptoten}.
  \item Bereken \(f'\), onderzoek het teken en bepaal \textbf{stijgen, dalen en extrema}.
  \item Bereken \(f''\), onderzoek het teken en bepaal \textbf{kromming en buigpunten}.
  \item Breng alle informatie samen in een \textbf{verloopschema}.
  \item Maak ten slotte een \textbf{grafiekschets} die met alle gegevens overeenstemt.
\end{enumerate}

\begin{opmerking}[Niet elk functietype vraagt dezelfde accenten]
Bij een rationale functie spelen verticale en eventueel horizontale of
schuine asymptoten vaak een belangrijke rol.

Bij een exponentiële functie zijn domein en continuïteit meestal
eenvoudig, maar gedrag op oneindig en groeisnelheid zijn essentieel.

Bij een logaritmische functie moet het domein zorgvuldig worden bewaakt.

Bij een goniometrische functie zijn periodiciteit en symmetrie vaak
belangrijker dan gedrag voor \(x\to\pm\infty\).
\end{opmerking}

% ============================================================
% 6. VOLLEDIG RATIONAAL VOORBEELD
% ============================================================

\FVDsection{Volledig voorbeeld: een rationale functie}

Onderzoek volledig

\[
f(x)=\frac{x^3}{x^2-1}.
\]

Dit voorbeeld sluit aan bij het oudere cursusmateriaal, maar wordt hier
systematisch en correct opgebouwd.

\FVDsubsection{Domein en continuïteit}

De noemer mag niet nul zijn:

\[
x^2-1=0
\Longleftrightarrow
x=\pm1.
\]

Dus

\[
\boxed{\operatorname{dom}f=\mathbb{R}\setminus\{-1,1\}.}
\]

Als rationale functie is \(f\) continu op elk interval van haar domein.

\FVDsubsection{Nulpunten, assen en teken}

De teller is nul voor

\[x^3=0\Longleftrightarrow x=0.\]

Dus de grafiek snijdt beide assen in

\[\boxed{(0,0)}.\]

Voor het teken onderzoeken we

\[
\frac{x^3}{(x-1)(x+1)}.
\]

De kritieke waarden zijn \(-1,0,1\). Daaruit volgt:

\[
\begin{array}{c|ccccccc}
x&(-\infty,-1)&-1&(-1,0)&0&(0,1)&1&(1,+\infty)\\
\hline
f(x)&-&\mid&+&0&-&\mid&+
\end{array}
\]

\FVDsubsection{Symmetrie}

\[
f(-x)=\frac{-x^3}{x^2-1}=-f(x).
\]

Dus \(f\) is oneven en de grafiek is puntsymmetrisch ten opzichte van
de oorsprong.

\FVDsubsection{Asymptoten}

Omdat de noemer nul wordt bij \(x=\pm1\), onderzoeken we eenzijdige
limieten. Dat levert verticale asymptoten

\[
\boxed{x=-1 \qquad\text{en}\qquad x=1.}
\]

De graad van de teller is precies één groter dan die van de noemer.
We voeren een veeltermdeling uit:

\[
\frac{x^3}{x^2-1}
=
x+\frac{x}{x^2-1}.
\]

Omdat

\[
\lim_{x\to\pm\infty}\frac{x}{x^2-1}=0,
\]

heeft de grafiek de schuine asymptoot

\[
\boxed{y=x.}
\]

De ligging ten opzichte van deze asymptoot wordt bepaald door

\[
f(x)-x=\frac{x}{x^2-1}.
\]

\FVDsubsection{Eerste afgeleide: stijgen, dalen en extrema}

\[
f'(x)
=
\frac{3x^2(x^2-1)-x^3(2x)}{(x^2-1)^2}
=
\frac{x^2(x^2-3)}{(x^2-1)^2}.
\]

De noemer is positief op het domein, dus het teken wordt bepaald door

\[x^2(x^2-3).\]

De nulpunten van \(f'\) zijn

\[x=-\sqrt3,\quad x=0,\quad x=\sqrt3.\]

Samen met de domeingrenzen \(x=\pm1\) krijgen we het tekenverloop:

\[
\begin{array}{c|ccccccccccc}
x & (-\infty,-\sqrt3) & -\sqrt3 & (-\sqrt3,-1) & -1 &
(-1,0) & 0 & (0,1) & 1 & (1,\sqrt3) & \sqrt3 & (\sqrt3,+\infty)\\
\hline
f'(x)&+&0&-&\mid&-&0&-&\mid&-&0&+
\end{array}
\]

Dus:

\[
x=-\sqrt3 \quad\text{geeft een lokaal maximum},
\]

\[
x=\sqrt3 \quad\text{geeft een lokaal minimum}.
\]

Bij \(x=0\) verandert het teken van \(f'\) niet, dus daar ligt geen
extremum.

De functiewaarden zijn

\[
f(-\sqrt3)=\frac{-3\sqrt3}{2},
\qquad
f(\sqrt3)=\frac{3\sqrt3}{2}.
\]

\FVDsubsection{Tweede afgeleide: kromming en buigpunten}

Uit

\[
f'(x)=\frac{x^2(x^2-3)}{(x^2-1)^2}
\]

volgt na vereenvoudiging

\[
\boxed{f''(x)=\frac{2x(x^2+3)}{(x^2-1)^3}.}
\]

Omdat \(x^2+3>0\), wordt het teken van \(f''\) bepaald door

\[
\frac{x}{(x^2-1)^3}.
\]

De mogelijke scheidingswaarden zijn \(-1,0,1\). We krijgen:

\[
\begin{array}{c|ccccccc}
x&(-\infty,-1)&-1&(-1,0)&0&(0,1)&1&(1,+\infty)\\
\hline
f''(x)&-&\mid&+&0&-&\mid&+
\end{array}
\]

Bij \(x=0\) verandert het teken van \(f''\), dus

\[
\boxed{(0,0)}
\]

is een buigpunt. Omdat ook \(f'(0)=0\), is dit een stationair buigpunt.

\FVDsubsection{Synthese}

Een correcte grafiekschets moet nu tegelijk voldoen aan:

\begin{itemize}
  \item puntsymmetrie ten opzichte van de oorsprong;
  \item verticale asymptoten \(x=-1\) en \(x=1\);
  \item schuine asymptoot \(y=x\);
  \item lokaal maximum bij \(x=-\sqrt3\);
  \item lokaal minimum bij \(x=\sqrt3\);
  \item stationair buigpunt in \((0,0)\);
  \item de gevonden stijg-, daal- en krommingsintervallen.
\end{itemize}

Dit is de kern van een functieonderzoek: de grafiek is het
\textbf{gevolg} van de analytische informatie, niet van gokken op basis
van enkele losse punten.

% ============================================================
% 7. ANDERE FUNCTIETYPES
% ============================================================

\FVDsection{Functieonderzoek buiten rationale functies}

Het leerdoel gaat over functies in het algemeen. Daarom moet je de
werkwijze ook kunnen overdragen naar andere functietypes.

\FVDsubsection{Irrationale functie}

Beschouw

\[f(x)=x\sqrt{x+2}.\]

Het domein is

\[[-2,+\infty).\]

Met de product- en kettingregel:

\[
f'(x)
=
\sqrt{x+2}+\frac{x}{2\sqrt{x+2}}
=
\frac{3x+4}{2\sqrt{x+2}}.
\]

Voor \(x>-2\) is de noemer positief. Het teken van \(f'\) wordt dus
bepaald door \(3x+4\). Het enige stationaire punt ligt bij

\[x=-\frac43.\]

Daar verandert \(f'\) van negatief naar positief, zodat \(f\) daar een
lokaal én absoluut minimum bereikt.

Hier zie je hoe domeinvoorwaarden en afgeleiden samen het verloop bepalen.

\FVDsubsection{Exponentiële functie}

Beschouw

\[f(x)=x e^{-x}.\]

Het domein is \(\mathbb{R}\). Met de productregel:

\[
f'(x)=e^{-x}(1-x).
\]

Omdat \(e^{-x}>0\), wordt het teken volledig bepaald door \(1-x\).
Dus \(f\) stijgt voor \(x<1\), daalt voor \(x>1\) en heeft een maximum
bij

\[
\boxed{\left(1,\frac1e\right)}.
\]

De tweede afgeleide is

\[
f''(x)=e^{-x}(x-2).
\]

Dus bij \(x=2\) verandert de kromming en ligt een buigpunt.

\FVDsubsection{Logaritmische functie}

Beschouw

\[f(x)=\frac{\ln x}{x},\qquad x>0.\]

De afgeleide is

\[
f'(x)=\frac{1-\ln x}{x^2}.
\]

Omdat \(x^2>0\), wordt het teken bepaald door \(1-\ln x\). Dus

\[
f'(x)=0
\Longleftrightarrow
\ln x=1
\Longleftrightarrow
x=e.
\]

De functie stijgt op \((0,e)\), daalt op \((e,+\infty)\) en heeft een
maximum in

\[
\boxed{\left(e,\frac1e\right)}.
\]

\FVDsubsection{Goniometrische functie}

Beschouw

\[f(x)=\sin x+\frac12\sin(2x).\]

Omdat de functie periodiek is, volstaat het vaak om één periode te
onderzoeken.

\[
f'(x)=\cos x+\cos(2x).
\]

Nu wordt het functieonderzoek gekoppeld aan goniometrische vergelijkingen:
de stationaire punten vind je door

\[
\cos x+\cos(2x)=0
\]

op het gekozen interval op te lossen.

Hier is dus niet alleen de afgeleide techniek belangrijk, maar ook de
keuze van een geschikte algebraïsche of goniometrische methode.

% ============================================================
% 8. EXTREMUMPROBLEMEN
% ============================================================

\FVDsection{Extremumproblemen: van context naar functie}

Bij een extremumprobleem is de functie meestal niet gegeven. Je moet ze
zelf opstellen uit de context.

Dat maakt deze problemen fundamenteel anders dan een standaardvraag
``bepaal de extrema van \(f\)''.

\begin{opmerking}[Modelleerstrategie]
Een bruikbare aanpak is:

\begin{enumerate}
  \item \textbf{Kies de variabele.}
  \item \textbf{Vertaal de voorwaarden} uit de context naar vergelijkingen.
  \item \textbf{Stel de doelfunctie op}: welke grootheid moet maximaal of minimaal worden?
  \item \textbf{Bepaal het betekenisvolle domein}. Contextvoorwaarden zijn hier essentieel.
  \item Bereken \(f'\) en bepaal kandidaat-extrema.
  \item \textbf{Controleer} dat het gevonden punt werkelijk het gevraagde extremum geeft.
  \item \textbf{Interpreteer} het resultaat in de oorspronkelijke context, met correcte eenheden.
\end{enumerate}
\end{opmerking}

\begin{voorbeeld}[Maximale oppervlakte]
Een rechthoek heeft een omtrek van \(40\,\mathrm{m}\).
Welke afmetingen geven de grootste oppervlakte?

Noem de lengte \(x\) en de breedte \(y\). Dan

\[
2x+2y=40
\Longrightarrow
y=20-x.
\]

De oppervlakte is

\[
A(x)=x(20-x)=20x-x^2.
\]

Uit de context volgt

\[0<x<20.\]

Nu

\[
A'(x)=20-2x.
\]

Dus

\[
A'(x)=0
\Longleftrightarrow
x=10.
\]

Verder is

\[A''(x)=-2<0,\]

dus \(x=10\) geeft een maximum.

De breedte is dan ook \(10\). De optimale rechthoek is dus een vierkant:

\[
\boxed{10\,\mathrm{m}\times10\,\mathrm{m}}.
\]

De maximale oppervlakte is

\[
\boxed{100\,\mathrm{m}^2}.
\]
\end{voorbeeld}

\begin{voorbeeld}[Minimale materiaalkost: open doos]
Uit een vierkant stuk karton met zijde \(30\,\mathrm{cm}\) worden in
de vier hoeken gelijke vierkantjes met zijde \(x\) weggeknipt. Daarna
worden de randen omhoog gevouwen.

Het volume van de open doos is

\[
V(x)=x(30-2x)^2.
\]

Uit de context volgt

\[0<x<15.\]

Als we het volume willen maximaliseren, onderzoeken we

\[
V'(x)
=
(30-2x)^2-4x(30-2x).
\]

Factoriseren geeft

\[
V'(x)=(30-2x)(30-6x).
\]

Binnen het betekenisvolle domein is

\[x=5\]

de relevante stationaire waarde. Een tekenonderzoek toont dat \(V'\)
daar van positief naar negatief verandert.

Dus het volume is maximaal voor

\[
\boxed{x=5\,\mathrm{cm}}.
\]

De doos heeft dan afmetingen

\[
20\,\mathrm{cm}\times20\,\mathrm{cm}\times5\,\mathrm{cm}
\]

en maximaal volume

\[
\boxed{2000\,\mathrm{cm}^3}.
\]
\end{voorbeeld}

\begin{voorbeeld}[Wetenschappelijke context: maximale vermogensoverdracht]
Een eenvoudige modelvorm voor het vermogen dat aan een belasting
\(R>0\) wordt geleverd is

\[
P(R)=\frac{E^2R}{(R+r)^2},
\]

waar \(E>0\) en \(r>0\) constanten zijn.

We zoeken voor welke \(R\) het vermogen maximaal is.

Omdat \(E^2>0\), volstaat het het verloop van

\[
p(R)=\frac{R}{(R+r)^2}
\]

te onderzoeken.

\[
p'(R)
=
\frac{(R+r)^2-2R(R+r)}{(R+r)^4}
=
\frac{r-R}{(R+r)^3}.
\]

Voor \(R>0\) is de noemer positief. Dus \(p'\) is positief voor \(R<r\)
en negatief voor \(R>r\).

Het vermogen is daarom maximaal wanneer

\[
\boxed{R=r.}
\]

Dit is een voorbeeld waarin een afgeleide rechtstreeks leidt tot een
technisch ontwerpcriterium.
\end{voorbeeld}

% ============================================================
% 9. ICT
% ============================================================

\FVDsection{ICT als controle- en onderzoeksinstrument}

Een grafiekprogramma of CAS kan zeer nuttig zijn bij functieonderzoek,
maar het vervangt de analyse niet.

Gebruik ICT bijvoorbeeld om:

\begin{itemize}
  \item een analytisch gevonden grafiekschets te controleren;
  \item vermoedelijke extrema of buigpunten te lokaliseren vóór je ze exact berekent;
  \item het effect van een parameter te onderzoeken;
  \item verschillende modellen te vergelijken;
  \item numerieke resultaten te controleren wanneer een exacte oplossing niet haalbaar is.
\end{itemize}

\begin{opmerking}[Analytisch eerst, digitaal bewust]
Een grafiekvenster kan belangrijke kenmerken verbergen of vervormen.
Een verticale asymptoot kan bijvoorbeeld als een bijna verticale lijn
verschijnen, en een extremum kan buiten het gekozen venster liggen.

Gebruik ICT daarom als \textbf{onderzoeks- en controlemiddel}, niet als
vervanging voor domeinonderzoek, limieten en tekenanalyse.
\end{opmerking}

% ============================================================
% 10. KORTE CONTROLE
% ============================================================

\FVDsection{Korte controle}

\begin{oefening}[Rolle en Lagrange]{\oefB\ESS}
\begin{question}
Onderzoek of Rolle toepasbaar is op
\[f(x)=\cos x\]
op \([0,2\pi]\). Bepaal alle waarden \(c\in(0,2\pi)\) waarvoor de
conclusie van Rolle geldt.

\begin{oplossing}
De functie is continu op \([0,2\pi]\), afleidbaar op \((0,2\pi)\) en
\[
f(0)=1=f(2\pi).
\]
Rolle is dus toepasbaar. Omdat
\[
f'(x)=-\sin x,
\]
moet
\[
\sin c=0.
\]
Binnen \((0,2\pi)\) geeft dit
\[
\boxed{c=\pi}.
\]
\end{oplossing}
\end{question}

\begin{question}
Toon met Lagrange dat er voor \(f(x)=\ln x\) op \([1,e]\) een
\(c\in(1,e)\) bestaat waarvoor
\[
f'(c)=\frac1{e-1}.
\]
Bepaal \(c\).

\begin{oplossing}
\(\ln x\) is continu op \([1,e]\) en afleidbaar op \((1,e)\).
Lagrange geeft
\[
f'(c)=\frac{\ln e-\ln1}{e-1}=\frac1{e-1}.
\]
Omdat
\[
f'(c)=\frac1c,
\]
volgt
\[
\boxed{c=e-1}.
\]
\end{oplossing}
\end{question}
\end{oefening}

\begin{oefening}[Van \(f'\) en \(f''\) naar conclusies]{\oefB\ESS}
Gegeven is
\[
f'(x)=\frac{(x-1)(x+2)}{(x-3)^2}
\]
voor een functie met \(x=3\) buiten het domein.

\begin{question}
Bepaal de intervallen waarop \(f\) stijgt en daalt.

\begin{oplossing}
De noemer is positief in het domein. Het teken wordt dus bepaald door
\[
(x-1)(x+2).
\]
Daarom is \(f'\) positief voor \(x<-2\) en \(x>1\), en negatief voor
\(-2<x<1\), met opsplitsing bij \(x=3\).

Dus \(f\) stijgt op
\[
(-\infty,-2),\qquad(1,3),\qquad(3,+\infty)
\]
en daalt op
\[
(-2,1).
\]
\end{oplossing}
\end{question}

\begin{question}
Welke stationaire waarden leveren een lokaal extremum?

\begin{oplossing}
Bij \(x=-2\) verandert \(f'\) van positief naar negatief: lokaal maximum.

Bij \(x=1\) verandert \(f'\) van negatief naar positief: lokaal minimum.
\end{oplossing}
\end{question}
\end{oefening}

\begin{oefening}[Analyseer, niet alleen bereken]{\oefU\ESS}
Voor
\[
f(x)=x e^{-x}
\]
berekende een leerling
\[
f'(x)=e^{-x}(1-x).
\]

Leg zonder een grafiekprogramma te gebruiken uit waarom de functie exact
één lokaal extremum heeft, welk type extremum dit is en waarom het zelfs
een absoluut maximum is.

\begin{oplossing}
Omdat \(e^{-x}>0\) voor alle \(x\), wordt het teken van \(f'\) volledig
bepaald door \(1-x\).

Voor \(x<1\) is \(f'(x)>0\), dus \(f\) stijgt. Voor \(x>1\) is
\(f'(x)<0\), dus \(f\) daalt. Bij \(x=1\) verandert het teken van positief
naar negatief, zodat daar een lokaal maximum ligt.

Verder geldt
\[
\lim_{x\to+\infty}xe^{-x}=0
\]
en
\[
\lim_{x\to-\infty}xe^{-x}=-\infty.
\]
Omdat de functie eerst stijgt tot \(x=1\) en daarna daalt, is het maximum
ook absoluut. De maximale functiewaarde is
\[
\boxed{f(1)=\frac1e}.
\]
\end{oplossing}
\end{oefening}

\begin{opmerking}[Wat moet je minimaal beheersen?]
Om het officiële leerplandoel op het vereiste niveau te bereiken, moet je
meer kunnen dan afgeleiden berekenen.

Je moet zeker de oefeningen met \(\ESS\) beheersen. Daarin moet je:

\begin{itemize}
  \item voorwaarden van Rolle en Lagrange controleren;
  \item informatie uit \(f'\) en \(f''\) interpreteren;
  \item zelf kritieke waarden selecteren;
  \item verschillende onderdelen van een functieonderzoek combineren;
  \item een conclusie over extrema of kromming verantwoorden;
  \item en bij extremumproblemen de wiskundige functie uit een context kunnen opstellen.
\end{itemize}
\end{opmerking}

\end{document}
